\section{Análise}

Em um primeiro momento, o simbolismo construído pelo texto remete à liberdade:
um pássaro voando sobre a praia, em demanda de céus mais frios. Entretanto, na
linha \texttt{6}, essa imagem é completamente quebrada. Os mesmos símbolos
introduzidos anteriormente passam a adquirir outro significado. As asas abertas
em cruz já não pertencem a uma ave em pleno voo, mas a um animal morto. As
areias molhadas, que antes evocavam a imagem de uma praia, convertem-se no
úmido areal onde repousa o cadáver. Da mesma forma, o bico aberto deixa de ser
o de uma ave em movimento e passa a guardar um grito que jamais será emitido.

Schmidt emprega um forte paralelismo gradativo ao longo do poema, acrescentando
a cada verso um novo elemento que amplia ou redefine a imagem anterior. Isso
pode ser observado, por exemplo, na progressão das linhas \texttt{4} a
\texttt{5}, bem como das linhas \texttt{7} a \texttt{9}. Em vez de repetir a
mesma ideia, cada novo verso intensifica o clima de tragédia e conduz o leitor
a reinterpretar as imagens já apresentadas.

Essa construção é muito mais rica do que qualquer formulação direta. Um verso
simples como ``Chovia na hora em que o contemplei'' produz uma forte inversão
de expectativa. O efeito não decorre do emprego de um vocabulário erudito, mas
da cena que o poeta cuidadosamente constrói no imaginário do leitor antes desse
momento. Com poucas palavras, Schmidt consegue deslocar completamente a
atmosfera do poema, conduzindo-a da contemplação da liberdade para a
constatação da morte. Enquanto poeta amador, dificilmente conseguiria alcançar
tamanho grau de sofisticação em tão poucos versos.
